% questions frame (shared by direction-slides and talk-slides)
\begin{frame}{The question: can statistical physics describe an agent swarm?}
\begin{columns}[T,onlytextwidth]
\begin{column}{0.36\textwidth}
\vspace{0.6em}
{\large\bfseries Goal}\par\vspace{0.4em}
A quantitative physics of LLM swarms that an operator can use: measured constants, fitted models with error bars, and laws that survive the four impostors.\par\vspace{0.8em}
{\large\bfseries Unit}\par\vspace{0.4em}
One goal period = one point on a phase diagram. Fit inside a period; compare periods by their parameters.
\end{column}
\begin{column}{0.60\textwidth}
\footnotesize
\renewcommand{\arraystretch}{1.45}
\begin{tabular}{@{}l >{\raggedright\arraybackslash}p{4.6cm} l@{}}
\textbf{Q1} & What couples agents? & \st{a}\\
\textbf{Q2} & What is field and what is coupling? & \st{a}\\
\textbf{Q3} & Is there collective order beyond fields? & \st{p}\\
\textbf{Q4} & Where does information live, and what is it worth? & \st{p}\\
\textbf{Q5} & What can an operator do? & \st{p}\\
\textbf{Q6} & Thermodynamics and selection & \st{o}\\
\textbf{Q7} & What transfers outside the village? & \st{o}\\
\end{tabular}\par\vspace{0.5em}
{\scriptsize\color{muted}``Answered'' means a result with high credence, not yet confirmed on reserved data.}
\end{column}
\end{columns}
\end{frame}
